\chapter*{Paisley}\stepcounter{chapter}\phantomsection\addcontentsline{toc}{chapter}{Paisley}
\DndDropCapLine{P}{\entryfont aisley sits on the ragged edge of the swamps like a settlement unsure whether it wishes to endure or surrender. The hamlet is small - no more than a few dozen structures - but visually divided against itself.}

{\entryfont On the western side stand the old homes: squat timber cottages sunken into damp earth, beams warped and blackened by moisture. Roofs sag. Fences lean. Smoke crawls from crooked chimneys. Elderly villagers linger outside at all hours, eyes fixed on the marsh, speaking in low tones of the "Living Death" and things that move where fog grows thickest.

Across a muddy lane stand new constructions: fresh-cut timber frames, straight fences, clean thatch not yet weathered by rot. Younger settlers dig drainage trenches and lay foundations with confident energy. They do not whisper about curses. They laugh at such talk. To them, the swamp is untapped potential - rich soil, trade routes, opportunity.

The tension is palpable but restrained. The generations coexist, but their world views are fundamentally incompatible.

Subtle remnants of older corruption remain in the land. The soil is rich - but overly so, almost unnaturally fertile. Pools of water sometimes reflect shapes not present above them. Livestock occasionally shy at empty air.

Paisley feels like a place caught between memory and ambition.}
\begin{DndReadAloud}
	The road dissolves into churned mud as you crest a low rise and see Paisley spread before you. The air smells of wet soil and fresh-cut wood - and beneath it all, the stagnant breath of the swamp. To one side, old cottages slump toward the earth, wood darkened by years of damp. Grey-haired villagers stand watchfully outside, gazes fixed on the distant swamps as if expecting something to stir.

	"Mark me", an old farmer mutters, gripping his hoe like a spear. "Nothing good ever came from that mire."

	A young woman a few paces away drives a stake into freshly turned soil, the timber frame of her future home rising behind her.

	"Nothing good?" she scoffs without looking at him. "It's mud and water, father. That's all it's ever been. You and the others just like scaring yourselves."

	"I've seen what crawls in fog."

	"You've seen shadows", she shoots back. "And every shadow doesn't mean a monster."

	A few other young settlers laugh, hoisting beams into place. One calls out, "If the swamp wanted us gone, it's had centuries to try."
\end{DndReadAloud}
\section*{Swamps of Paisley}\stepcounter{section}\phantomsection\addcontentsline{toc}{section}{Swamps of Paisley}
\subsection*{The Altar of Evil}
{\entryfont At first, the swamp appears ecologically typical: shallow brackish pools, peat-heavy soil, cypress-like trees with exposed root systems, curtains of moss, and the constant drone of insects. Amphibians scatter at intrusion. Waterfowl lift lazily from reeds. Small reptiles slide into murk at the party's approach.

However, the deeper the party proceeds, the biome shifts subtly but unmistakably.
\begin{itemize}
	\item Natural animal calls diminish in frequency.
	\item The insect drone becomes intermittent.
	\item Tracks become scarce.
	\item The water grows darker, more stagnant.
	\item The air grows still — unnaturally so.
\end{itemize}
The silence does not feel peaceful. It feels withheld.

Occasionally, the party is ambushed by Twig Blights or Needle Blights emerging from dense root clusters or reed beds. These attacks are sudden but tactically unsophisticated. The blights fight without coordination and crumble quickly. Their presence suggests corruption in the flora rather than organized hostility.

The further inward the party travels, the more plant life appears twisted: bark splitting in unnatural spirals, thorned vines growing in circular patterns, small patches of vegetation blackened as though frostbitten despite the damp warmth.

Eventually, the forest canopy thins abruptly. The party emerges into a circular clearing.

At its center stands an ancient stone tablet or altar — weathered, rectangular, and partially sunken into peat. The stone is carved with shallow grooves worn nearly smooth by time.

Surrounding it — and scattered across the clearing floor — lie skeletal remains.

Closer examination reveals:
\begin{itemize}
	\item Long, slender equine skulls.
	\item Rib cages proportioned for powerful lungs.
	\item Hoof structures consistent with large ungulates.
	\item Most disturbingly: several skulls bear a single central cranial cavity where a horn once anchored.
\end{itemize}
One intact unicorn horn lies partially buried near the altar, its surface dulled and cracked but unmistakable.

The number of skeletons suggests this was not a single event. It was repeated. The clearing feels like something concluded here long ago — yet not entirely resolved.}

\begin{DndReadAloud}
	The swamp thins without warning.

	Trees recede, their twisted roots giving way to a circular clearing where the air hangs heavy and unmoving.

	At its center rises a weathered stone slab, half-sunken into dark peat. Time has worn its edges smooth.

	Bones lie everywhere.

	Long-limbed skeletons sprawl across the clearing floor and drape over the stone itself, ribs arched toward the grey 	sky. Their skulls are unmistakably equine — yet each bears a single hollow at its brow.

	Near the altar, half-buried in the mud, rests a spiral horn.

	The swamp is silent.

	Completely silent.

	Not even insects dare intrude here.
\end{DndReadAloud}